\startchapter[title={Verificación de las soluciones}, reference=Anexo_Verificacion]

Este anexo detalla el verificador de soluciones descripto en la
Sección~\in[sec:validacion] y, al final, cómo se detecta la espera
voluntaria en las soluciones del Split.

%----------------------------------------------------------------------------------------
\startsection[title={Chequeos directos}]

Antes de usar el resolvedor, la solución se recalcula a partir de los datos
de la instancia y se verifica que:
\startitemize
	\item cada lote tenga un único {\em picker} y que ese {\em picker}
	exista;
	\item las rutas usen solo ubicaciones y arcos del grafo, no repitan
	ubicaciones y cubran todas las ubicaciones requeridas de su lote;
	\item cada lote cumpla la cadena de frío con el {\em picker} asignado,
	con la misma tolerancia que los modelos lineales ($\varepsilon = 10^{-9}$);
	\item los tiempos de ejecución y los instantes de finalización informados
	coincidan con los recalculados;
	\item el objetivo informado coincida con el que da la fórmula del modelo
	monolítico (Ecuación~\in[eq:mono_objetivo]).
\stopitemize
Cada chequeo que falla se registra con un mensaje que indica la restricción y
el lote involucrados.

\stopsection

%----------------------------------------------------------------------------------------
\startsection[title={Cómo se fija una solución en el modelo}]

El modelo monolítico se compila con ZIMPL a un archivo en formato LP, en el
que cada variable tiene un nombre formado por la variable y sus índices. El
verificador lee ese archivo y fija las variables de la solución por nombre,
igualando sus cotas inferior y superior al valor correspondiente, sin
modificar ninguna restricción. En el nivel~1 se fijan $z_{ijb}$ (1 en los
arcos de la ruta del lote y 0 en el resto), $x_{pb}$ (1 para el
{\em picker} asignado y 0 para los demás) e $y_{bk}$ (según el orden de los
lotes en la secuencia de cada {\em picker}). En el nivel~2 se fijan, además,
los instantes de finalización $c_{pb}$ y $c_b$ informados. Las demás
variables, como el flujo $q_{ijb}$, el tiempo de ruta $r_b$, el tiempo de
ejecución $u_{pb}$ y el adelanto y la demora de cada orden, quedan libres y
las completa el resolvedor. Si el nivel~1 resulta infactible, el verificador
busca qué lote lo provoca, para que el reporte indique dónde está la
violación.

\stopsection

%----------------------------------------------------------------------------------------
\startsection[title={Ajustes del verificador}]

Para que la verificación pueda aplicarse a todas las instancias, el
verificador hace tres ajustes, que quedan anotados en su reporte:
\startitemize[n]
	\item {\bf Lotes vacíos.} Quita de la instancia los lotes sin pedidos que
	la solución no usa. Los modelos lineales rutean, asignan y cobran todos
	los lotes (Sección~\in[sec:caracteristicas_formulacion]), mientras que las
	heurísticas ignoran los vacíos.
	\item {\bf Subgrafo de las rutas.} Cuando todas las rutas están fijas,
	resuelve sobre el subgrafo formado por los arcos de esas rutas, el
	depósito y el punto de entrega. Es exactamente equivalente, porque un arco
	que no usa ninguna ruta vale 0 en todas, y resultó necesario en el almacén
	de 100×100, donde ZIMPL se quedaba sin memoria al compilar el modelo
	completo. Como control, 248 soluciones dieron el mismo veredicto, los
	mismos estados y los mismos objetivos con y sin este recorte.
	\item {\bf Constantes M grandes.} Si las constantes $M$ o $M^{\text{seq}}$
	de la instancia son demasiado chicas para la solución, las reemplaza por
	valores válidos, porque una constante chica puede excluir soluciones
	factibles (Sección~\in[sec:mono_asignacion]) y el rechazo no se debería a
	una violación.
\stopitemize

\stopsection

%----------------------------------------------------------------------------------------
\startsection[title={Validación del verificador}]

Para comprobar que el verificador detecta errores, se introdujeron cuatro a
propósito en una solución chica: un arco que no existe en el grafo, dos lotes
del mismo {\em picker} superpuestos en el tiempo, una violación de la cadena
de frío y un objetivo mal informado. El verificador detectó los cuatro.

\stopsection

%----------------------------------------------------------------------------------------
\startsection[title={Pseudocódigo}]

La Figura~\in[fig:verificador_pseudocodigo] resume el procedimiento que el
verificador aplica a cada solución. La solución $s = (\pi, a, \sigma)$ tiene
la forma de la Sección~\in[sec:hc_representacion], y la heurística informa
además el tiempo de ejecución $\hat{T}_b$ y el instante de finalización
$\hat{c}_b$ de cada lote, y el valor objetivo $\hat{f}$. Los valores sin
sombrero son los que recalcula el verificador.

\startplacefigure[title={Pseudocódigo del verificador de soluciones.}, reference=fig:verificador_pseudocodigo]
\startframedtext[width=broad, frame=on, align=flushleft]
\startlines
{\bf Entrada:} instancia; solución $s = (\pi, a, \sigma)$ con los valores informados $\hat{T}_b$, $\hat{c}_b$ y $\hat{f}$
{\bf Salida:} válida o con violaciones, y un registro $R$ con los motivos
\blank[small]
{\em Preparación}
1.\quad quitar de $\mathcal{B}$ los lotes sin órdenes que no aparecen en $s$
2.\quad $R \leftarrow \emptyset$;\; $\mathcal{F} \leftarrow \emptyset$ \hfill {\em lotes que se fijarán en el modelo}
{\em Chequeos directos}
3.\quad {\bf para cada} lote $b \in \mathcal{B}$ {\bf hacer}
4.\quad\qquad {\bf si} $b$ no tiene exactamente un {\em picker} en $s$ {\bf entonces}
\dontleavehmode\hphantom{00.}\quad\qquad\qquad registrar en $R$ y seguir con el próximo lote
5.\quad\qquad {\bf si} la ruta $(v_{\text{start}}, \pi_b, v_{\text{end}})$ usa ubicaciones o arcos inexistentes {\bf entonces}
\dontleavehmode\hphantom{00.}\quad\qquad\qquad registrar en $R$ y seguir con el próximo lote
6.\quad\qquad {\bf si} la ruta repite ubicaciones o no cubre las requeridas por $b$ {\bf entonces}
\dontleavehmode\hphantom{00.}\quad\qquad\qquad registrar en $R$
7.\quad\qquad $T_b \leftarrow F_b + r_b / V_{a(b)}$, con el tiempo de ruta $r_b$ calculado sobre $\pi_b$
8.\quad\qquad {\bf si} $T_b > T^{\text{cold}} + \varepsilon$ {\bf o} $T_b \neq \hat{T}_b$ {\bf entonces} registrar en $R$
9.\quad\qquad agregar $b$ a $\mathcal{F}$
10.\quad {\bf para cada} {\em picker} $p$ {\bf hacer}
11.\quad\qquad recorrer $\sigma_p$ desde el instante 0 y calcular cada $c_b$ acumulando los $T_b$
12.\quad\qquad {\bf si} algún $c_b \neq \hat{c}_b$ {\bf entonces} registrar en $R$
13.\quad calcular $f$ con la Ecuación~\in[eq:mono_objetivo] y los $c_b$
14.\quad {\bf si} $f \neq \hat{f}$ {\bf entonces} registrar en $R$
{\em Modelo}
15.\quad {\bf si} $\mathcal{F} = \mathcal{B}$ {\bf entonces} reducir el grafo a las ubicaciones y los arcos de las rutas
16.\quad {\bf si} $M$ o $M^{\text{seq}}$ son chicas para los tiempos de $s$ {\bf entonces}
\dontleavehmode\hphantom{00.}\quad\qquad reemplazarlas por valores válidos
17.\quad compilar el modelo monolítico con ZIMPL
{\em Nivel 1: decisiones fijas}
18.\quad fijar $z_{ijb}$ y $x_{pb}$ para los lotes de $\mathcal{F}$, e $y_{bk}$ según $\sigma$; resolver con SCIP
19.\quad {\bf si} el modelo es infactible {\bf entonces}
20.\quad\qquad liberar por grupos las decisiones fijadas y registrar en $R$
\dontleavehmode\hphantom{00.}\quad\qquad\qquad cuáles restablecen la factibilidad
21.\quad\qquad {\bf devolver} con violaciones, $R$
{\em Nivel 2: solución completa fija}
22.\quad fijar además $c_{pb}$ y $c_b$ en los valores $\hat{c}_b$; resolver con SCIP; sea $f^*$ su objetivo
23.\quad {\bf si} $R = \emptyset$, el nivel 2 tiene solución y $f^* = \hat{f}$ {\bf entonces devolver} válida
24.\quad {\bf devolver} con violaciones, $R$
\stoplines
\stopframedtext
\stopplacefigure

Las comparaciones entre valores informados y recalculados se hacen con una
tolerancia relativa de $10^{-6}$. Los lotes que no llegan a $\mathcal{F}$,
porque su ruta usa ubicaciones o arcos inexistentes, no se fijan: quedan
libres en los niveles de SCIP, y su violación ya está registrada. En el
diagnóstico del paso~20, las decisiones se liberan primero por tipo (rutas,
secuencia, y asignación junto con secuencia) y después lote por lote, para
indicar qué ruta provoca la infactibilidad. Si la verificación no puede
ejecutarse, por ejemplo porque el modelo no compila, el resultado se informa
como un error del verificador y no como un veredicto sobre la solución. Las
instancias en las que la heurística no encontró solución no tienen nada que
fijar y no se verifican.

% TODO: no estoy segura de esto en general, podemos ver si queda o no. quizas ya
% es mucho ruido
\stopsection

%----------------------------------------------------------------------------------------
\startsection[title={Detección de la espera voluntaria}, reference=sec:anexo_espera]

Para decidir si una instancia se compara también en objetivo
(Sección~\in[sec:comparacion]), se determina si la solución del Split usa
la espera voluntaria (Sección~\in[sec:espera_voluntaria]). Para eso se
reconstruye la agenda de cada {\em picker} en esa solución: un lote espera si
figura como terminado después del instante en que podría haber terminado, es
decir, la finalización del lote anterior del mismo {\em picker} (o cero, si
es el primero) más su tiempo de ejecución. La espera solo se tiene en cuenta
si el adelanto pesa en el objetivo ($W^{\text{pen}} W^{e} > 0$), porque solo
entonces puede bajar el objetivo.

\stopsection
\stopchapter
